% SPDX-License-Identifier: PMPL-1.0-or-later
% SPDX-FileCopyrightText: 2026 Jonathan D.A. Jewell
%
% Groove: Formally Verified System Composition via
% Dependent Types and Linear Resources
%
% Intended venue: arXiv preprint (cs.SE / cs.PL)
% Subsequently: ICFP, POPL, or ECOOP

\documentclass[11pt,a4paper]{article}
\usepackage[utf8]{inputenc}
\usepackage[T1]{fontenc}
\usepackage{amsmath,amssymb,amsthm}
\usepackage{listings}
\usepackage{hyperref}
\usepackage{booktabs}
\usepackage{graphicx}
\usepackage{xcolor}
\usepackage[margin=2.5cm]{geometry}

\definecolor{codegreen}{rgb}{0.25,0.49,0.31}
\definecolor{codegray}{rgb}{0.5,0.5,0.5}
\definecolor{codeblue}{rgb}{0.22,0.40,0.72}

\lstset{
  basicstyle=\ttfamily\small,
  keywordstyle=\color{codeblue}\bfseries,
  commentstyle=\color{codegray}\itshape,
  stringstyle=\color{codegreen},
  breaklines=true,
  frame=single,
  captionpos=b
}

\newtheorem{theorem}{Theorem}
\newtheorem{lemma}{Lemma}
\newtheorem{definition}{Definition}
\newtheorem{property}{Property}

\title{Groove: Formally Verified System Composition via\\
Dependent Types and Linear Resources}

\author{Jonathan D.A. Jewell\\
\texttt{j.d.a.jewell@open.ac.uk}\\
The Open University}

\date{March 2026}

\begin{document}

\maketitle

\begin{abstract}
We present Groove, a protocol for runtime discovery and type-safe composition
of independent software systems. Unlike service meshes, plugin architectures,
or message brokers, Groove enables structural composition where compatibility
is a compile-time property enforced by dependent types and connection handles
are linear resources that cannot leak. Systems discover each other via
lightweight HTTP endpoint probing, negotiate capabilities through structural
type matching, and compose --- or remain independent --- without configuration.
We formalise the protocol in Idris2, prove four safety properties (match
safety, resource linearity, capability authenticity, and composition
soundness), and describe a reference implementation spanning Gossamer (desktop
shell), Burble (voice platform), VeriSimDB (multi-modal database), and Vext
(integrity verification). We show that Groove occupies a previously unfilled
design space: zero-configuration system composition with formal correctness
guarantees.
\end{abstract}

\section{Introduction}

Modern software systems are composed of independent components that must
interoperate. Three dominant integration models exist:

\begin{enumerate}
\item \textbf{Hard-coded integration}: systems import each other's libraries.
  Changes in one break the other. Coupling is total.
\item \textbf{Plugin architectures}: a host application defines extension
  points. Plugins are subordinate to the host. The relationship is
  asymmetric.
\item \textbf{Service meshes}: systems communicate over the network via
  proxied traffic. Coupling is loose, but compatibility is a runtime
  property discovered through failures.
\end{enumerate}

Each model makes a trade-off between coupling strength and safety guarantees.
Hard-coded integration provides compile-time safety but maximal coupling.
Service meshes provide minimal coupling but no compile-time safety. Plugin
architectures sit between, with host-defined contracts that plugins may
violate at runtime.

We propose a fourth model: \textbf{structural composition}. In Groove,
systems declare what they offer and what they consume via typed capability
interfaces. If the structural types match, the systems compose. If not,
both continue functioning independently. Compatibility is not a runtime
property --- it is a compile-time proof.

\subsection{Contributions}

\begin{enumerate}
\item A protocol for zero-configuration runtime discovery of typed service
  capabilities via HTTP endpoint probing (Section~\ref{sec:discovery}).
\item A formalisation of capability matching as structural subtyping with
  dependent types in Idris2, providing compile-time composition proofs
  (Section~\ref{sec:types}).
\item A linear resource discipline for connection handles that makes resource
  leaks compile-time errors (Section~\ref{sec:linear}).
\item A provenance attestation mechanism that generates verifiable audit
  trails for all composition events (Section~\ref{sec:attestation}).
\item A modular architecture following an ABI/FFI/API triple that enables
  polyglot participation without compromising type safety
  (Section~\ref{sec:polyglot}).
\item A zero-copy, frame-level IPv4$\rightarrow$IPv6 proxy implemented in Zig
  that transparently translates legacy connections using \texttt{splice(2)},
  requiring no elevated privileges (Section~\ref{sec:transport}).
\item A reference implementation across four independent systems demonstrating
  the protocol's practical applicability (Section~\ref{sec:implementation}).
\end{enumerate}

\section{Background and Related Work}

\subsection{Service Discovery}

DNS-SD \cite{rfc6763} and mDNS \cite{rfc6762} provide service discovery on
local networks. Consul and etcd provide distributed service registries. These
systems discover \emph{that} services exist and \emph{where} they are, but
not \emph{whether they are compatible}. Groove extends discovery with
structural compatibility checking.

\subsection{Capability-Based Security}

The capability model \cite{dennis1966} restricts access to resources via
unforgeable tokens. Groove uses capabilities not for access control but for
\emph{composition safety}: a capability token proves that a system implements
a specific interface. The token is a type-level proof, not a runtime
credential.

\subsection{Session Types and Behavioural Types}

Session types \cite{honda1993} guarantee protocol conformance in
communication. Groove's capability interfaces are related but distinct:
session types describe \emph{how} two parties communicate (message
sequences), while Groove capabilities describe \emph{what} parties can do
together (structural compatibility). The two are complementary.

\subsection{Dependent Types for Systems Programming}

Idris2 \cite{brady2021} provides dependent types with quantitative type
theory (QTT), enabling linear resource tracking at the type level. F*
\cite{fstar} and Agda \cite{norell2007} provide similar expressive power
but lack Idris2's practical compilation to efficient code via its Chez
Scheme or RefC backends.

\section{Protocol Design}

\subsection{Discovery}
\label{sec:discovery}

Every Groove-aware system exposes an HTTP endpoint at
\texttt{/.well-known/groove} (following RFC~8615) that returns a
\emph{capability manifest} in A2ML format:

\begin{lstlisting}[caption={Groove manifest for a voice platform}]
@groove-manifest(version="0.1.0"):
  @system(id="burble", version="1.2.0")
  @offers:
    @capability(id="voice", version="2.1.0",
                interface="VoiceTransport")
    @capability(id="text",  version="1.0.0",
                interface="TextChannel")
  @end
  @consumes:
    @capability(id="integrity", version="1.0+",
                interface="IntegrityVerifier")
  @end
@end
\end{lstlisting}

Discovery proceeds by probing a configurable port range (default 6460--6500)
on the loopback interface. IPv6 (\texttt{::1}) is probed first, with IPv4
(\texttt{127.0.0.1}) as fallback, following the Happy Eyeballs algorithm
(RFC~8305). Each probe uses a 500ms timeout. Full discovery completes within
5 seconds.

\subsection{Structural Capability Matching}
\label{sec:types}

Capability matching is defined as structural subtyping on interface
definitions:

\begin{definition}[Capability Match]
A provider capability $P$ matches a consumer requirement $C$ iff:
\begin{enumerate}
\item The interface type of $P$ is a structural supertype of $C$'s
  required interface type (the provider offers at least what the consumer
  needs).
\item The version of $P$ satisfies $C$'s version constraint (SemVer).
\item The security requirements of $C$ are met by $P$'s declared
  security properties.
\end{enumerate}
\end{definition}

In Idris2, this is expressed as an auto-implicit proof obligation:

\begin{lstlisting}[language=Haskell,caption={Type-safe connection in Idris2}]
connect : (prov : System)
       -> (cons : System)
       -> {auto prf : CapMatch (offers prov)
                               (requires cons)}
       -> IO (1 _ : GrooveConn prov cons)
\end{lstlisting}

The \texttt{auto} keyword instructs the compiler to search for a proof
of \texttt{CapMatch} automatically. If no proof exists --- because the
capabilities are structurally incompatible --- the code \emph{does not
compile}. This is not a runtime check; it is a compile-time impossibility.

\subsection{Linear Resource Safety}
\label{sec:linear}

Connection handles are linear resources (quantity 1 in QTT):

\begin{lstlisting}[language=Haskell,caption={Linear connection handle}]
data GrooveConn : System -> System -> Type where
  MkGroove : (handle : Bits64)
          -> GrooveConn prov cons

disconnect : (1 _ : GrooveConn p c) -> IO ()
\end{lstlisting}

The \texttt{1} quantity annotation means the connection handle must be
consumed \emph{exactly once}. The compiler rejects programs that:

\begin{itemize}
\item Forget to disconnect (the handle is never consumed --- \emph{leak}).
\item Disconnect twice (the handle is consumed more than once ---
  \emph{double-free}).
\item Use the handle after disconnection (the handle has been consumed ---
  \emph{use-after-free}).
\end{itemize}

\begin{theorem}[Resource Safety]
In any well-typed Groove program, every connection handle is consumed
exactly once. No connection leaks, double-disconnects, or
use-after-disconnect are possible.
\end{theorem}

\begin{proof}
Follows directly from QTT's linearity discipline. A value of quantity 1
must appear exactly once in the program's normal form. The Idris2 type
checker enforces this structurally.
\end{proof}

\subsection{Provenance Attestation}
\label{sec:attestation}

Every Groove connection event generates a provenance record linked by
SHA-256 hash chain:

\begin{lstlisting}[caption={Provenance record structure}]
{
  "event": "groove:connected",
  "provider": "burble@1.2.0",
  "consumer": "gossamer@0.8.0",
  "capabilities": ["voice@2.1.0", "text@1.0.0"],
  "timestamp": "2026-03-23T14:30:00Z",
  "hash": "sha256:a1b2c3...",
  "prev_hash": "sha256:e5f6g7..."
}
\end{lstlisting}

When VeriSimDB is available (discovered via Groove itself), provenance
records are stored as octads with full 8-modality representation.

\subsection{Graceful Degradation}

\begin{property}[Bare-First]
Every Groove-aware system functions correctly with zero groove partners.
Groove capabilities are additive; no capability is required.
\end{property}

When a partner disconnects (or its heartbeat times out after 15 seconds),
the consumer's enhanced capabilities disappear and the system reverts to
its base functionality. This transition is not an error --- it is the
designed behaviour.

\section{Polyglot Participation}
\label{sec:polyglot}

Systems not written in Idris2 participate through a three-layer
architecture:

\begin{enumerate}
\item \textbf{ABI (Idris2)}: interface definitions with dependent type
  proofs. Defines \emph{what} the interface is.
\item \textbf{FFI (Zig)}: C-compatible implementation with zero runtime
  overhead. Implements \emph{how} to call it from any language.
\item \textbf{API (V)}: ergonomic language-specific binding generators.
  Provides \emph{pleasant} APIs for each target language.
\end{enumerate}

Any language that can call C functions participates at full safety. The
Idris2 proofs are compiled into the C ABI boundary; callers need not be
aware of the proof machinery.

\section{Formal Properties}

We prove four safety properties of the Groove protocol:

\begin{theorem}[Match Safety]
If a Groove connection is established between systems $A$ and $B$, then
the offered capabilities of $A$ structurally satisfy the consumed
requirements of $B$ (and vice versa). Incompatible connections cannot be
expressed in the type system.
\end{theorem}

\begin{theorem}[Resource Linearity]
Every connection handle in a well-typed Groove program is consumed exactly
once. No resource leaks, double-frees, or use-after-free are possible.
\end{theorem}

\begin{theorem}[Capability Authenticity]
A system declaring capability $C$ in its manifest must implement the
interface type associated with $C$. Phantom capabilities (declared but
unimplemented) are rejected by the type checker.
\end{theorem}

\begin{theorem}[Composition Soundness]
If systems $A$, $B$, and $C$ form a groove mesh where $A \leftrightarrow B$
and $B \leftrightarrow C$, the composed system $A \leftrightarrow B
\leftrightarrow C$ is well-typed. Composition preserves all safety
properties of individual connections.
\end{theorem}

\section{Implementation}
\label{sec:implementation}

We have implemented Groove across four independent systems:

\begin{table}[h]
\centering
\begin{tabular}{llll}
\toprule
System & Language & Capabilities Offered & Port \\
\midrule
Gossamer & Ephapax/Zig & Desktop shell, panels & 6470 \\
Burble & Elixir & Voice, text, presence, TTS & 6473 \\
VeriSimDB & Zig & 8-modality storage, drift & 6475 \\
Vext & Zig & Integrity, hash chains & 6480 \\
\bottomrule
\end{tabular}
\caption{Reference implementation systems}
\end{table}

Each system functions independently. Starting any combination produces
the expected composed behaviour:

\begin{itemize}
\item Gossamer + Burble: desktop app with embedded voice UI
\item Burble + Vext: voice with cryptographic message integrity
\item Gossamer + VeriSimDB: persistent state with drift detection
\item All four: full groove mesh with verified voice and persistent state
\end{itemize}

\subsection{FLI: Frontend Library Interface}

For GUI panels within Gossamer, we implement a panel-level composition
system called FLI (Frontend-user Library Interface). Panels declare
\emph{clade traits} (e.g., \texttt{fli-gauge}, \texttt{fli-editable})
and the FLI provisioner loads matching JavaScript modules on-demand. This
follows the same bare-first, capability-layered pattern as Groove itself.

\subsection{Browser Extension Harness}

We describe a browser extension that acts as a Groove harness, enabling
web applications to discover and compose with local groove-aware services.
The extension probes \texttt{[::1]} (IPv6) and \texttt{127.0.0.1} (IPv4)
on the standard port range, respecting dual-stack requirements and CORS
constraints. This provides a type-safe bridge between the browser security
sandbox and native desktop services.

\section{Transport: IPv6-First with Frame-Level Translation}
\label{sec:transport}

\subsection{IPv6 as Default Transport}

Groove specifies IPv6-only binding (\texttt{[::1]:PORT}) as the recommended
configuration for localhost discovery. Every operating system capable of
running the Groove toolchain (Idris2, Zig) has a functional IPv6 loopback
interface. IPv4 is treated as a legacy transport with an explicit sunset
path.

For LAN discovery beyond localhost, Groove uses mDNS/DNS-SD over IPv6
multicast (\texttt{ff02::fb}), leveraging link-local addresses (\texttt{fe80::})
that require no DHCP configuration. LAN discovery is explicitly opt-in; the
default scope is localhost-only.

\subsection{Frame-Level IPv4$\rightarrow$IPv6 Proxy}

For systems that must accept IPv4 connections during a transition period,
we introduce a frame-level proxy that operates at the TCP stream layer
(Layer~4) rather than the HTTP application layer (Layer~7). This provides
transparent, zero-copy IPv4$\rightarrow$IPv6 translation.

The proxy accepts IPv4 TCP connections and opens a corresponding IPv6 TCP
connection to the Groove service on \texttt{[::1]:PORT}. Bytes are spliced
bidirectionally using the kernel's \texttt{splice(2)} syscall, which
transfers data between socket file descriptors without copying to userspace.

\begin{lstlisting}[caption={Frame-level proxy data path}]
Client (IPv4) --> [accept] --> Frame Proxy --> [splice] --> Service (IPv6)
Client (IPv4) <-- [splice] <-- Frame Proxy <-- [splice] <-- Service (IPv6)
\end{lstlisting}

This approach has several advantages over HTTP-level upgrade directives:

\begin{enumerate}
\item \textbf{Zero-copy}: On Linux, \texttt{splice(2)} moves data between
  sockets via kernel pipe buffers. No bytes are copied to userspace.
\item \textbf{Transparent}: The client is unaware of the translation. No
  reconnection, no upgrade headers, no protocol awareness needed.
\item \textbf{Protocol-agnostic}: Works for any protocol inside the TCP
  stream --- HTTP, WebSocket, gRPC, binary frames, or raw bytes. The proxy
  does not parse content.
\item \textbf{No elevated privileges}: Runs in userspace without kernel
  modules, \texttt{nftables} rules, or eBPF programs.
\end{enumerate}

The proxy is implemented in approximately 60 lines of Zig, with
\texttt{splice(2)} on Linux and a 4KB userspace buffer fallback on
macOS and Windows.

\begin{property}[Transport Transparency]
The frame-level proxy preserves all Groove protocol properties. From the
perspective of the Groove service, all connections arrive on IPv6. From
the perspective of the client, the connection succeeds on IPv4 with no
visible translation.
\end{property}

For browser extensions, which cannot perform raw TCP operations, the
protocol falls back to HTTP-level upgrade headers:
\texttt{Groove-IPv6-Preferred} and \texttt{Groove-Upgrade-To} directives
instruct the extension to reconnect on \texttt{[::1]} for subsequent
requests. This is analogous to HSTS (HTTP Strict Transport Security)
applied to the network layer rather than the encryption layer.

\subsection{CORS and the Browser Boundary}

All Groove HTTP endpoints include CORS headers to enable browser extension
access. Control-plane traffic uses HTTP/1.1; data-plane traffic may use
any transport (WebSocket, gRPC, WebRTC, Unix sockets) as declared in the
capability interface definition. The Idris2 ABI layer provides type-safe
proofs that cross the browser--native boundary, ensuring structural
compatibility between the browser security sandbox and local services.

\section{Evaluation}

\subsection{Discovery Latency}

Local discovery across 40 ports completes in under 2 seconds on commodity
hardware (dual-stack, 500ms timeout per probe, parallel probing). Discovery
of a running groove partner typically completes in under 100ms (first port
hit).

\subsection{Connection Overhead}

Groove connection establishment (manifest exchange + structural type check)
adds approximately 5ms over a raw TCP connection. Heartbeat traffic
consumes negligible bandwidth (one HTTP 204 response every 5 seconds per
connection).

\subsection{Proof Compilation}

The Idris2 proofs for match safety, linearity, authenticity, and
composition soundness compile in under 3 seconds on a 4-core machine.
Proof terms are erased at runtime; there is zero runtime overhead from
the formal verification.

\section{Discussion}

\subsection{Comparison with Existing Approaches}

\begin{table}[h]
\centering
\small
\begin{tabular}{lcccc}
\toprule
Property & Hard-coded & Plugins & Service Mesh & \textbf{Groove} \\
\midrule
Compile-time safety & \checkmark & Partial & $\times$ & \checkmark \\
Zero configuration & $\times$ & $\times$ & $\times$ & \checkmark \\
Peer relationship & $\times$ & $\times$ & \checkmark & \checkmark \\
Formal proofs & $\times$ & $\times$ & $\times$ & \checkmark \\
Resource safety & Partial & $\times$ & $\times$ & \checkmark \\
Graceful degradation & $\times$ & Partial & \checkmark & \checkmark \\
Provenance trail & $\times$ & $\times$ & $\times$ & \checkmark \\
\bottomrule
\end{tabular}
\caption{Comparison of integration models}
\end{table}

Groove is the only model that provides all seven properties simultaneously.

\subsection{Limitations}

\begin{enumerate}
\item The Idris2 type proofs are available only to systems that use the
  Idris2 ABI directly. Polyglot participants receive runtime validation
  rather than compile-time proofs.
\item Local discovery via port probing is inherently limited to the same
  machine. LAN and WAN discovery require additional infrastructure (mDNS,
  DNS-SD, or a registry).
\item The protocol does not currently address capability versioning
  migrations (deprecating a capability version across a mesh).
\end{enumerate}

\section{Conclusion}

Groove demonstrates that system composition can be both zero-configuration
\emph{and} formally verified. By combining dependent types for structural
matching, linear resources for connection safety, and lightweight HTTP
probing for discovery, Groove fills a design space that no existing
integration model occupies.

The protocol is intentionally simple: expose a manifest, probe nearby,
match structurally, compose or don't. The formal machinery is invisible
to users --- they see services that find each other and work together.
The types are the engine, not the paint job.

\bibliographystyle{plain}
\begin{thebibliography}{10}

\bibitem{brady2021}
E.~Brady.
\newblock \emph{Idris 2: Quantitative Type Theory in Practice}.
\newblock In ECOOP, 2021.

\bibitem{dennis1966}
J.~Dennis and E.~Van~Horn.
\newblock Programming semantics for multiprogrammed computations.
\newblock \emph{Communications of the ACM}, 9(3):143--155, 1966.

\bibitem{fstar}
N.~Swamy, C.~Hri\c{t}cu, C.~Keller, et al.
\newblock Dependent types and multi-monadic effects in {F*}.
\newblock In POPL, 2016.

\bibitem{honda1993}
K.~Honda.
\newblock Types for dyadic interaction.
\newblock In CONCUR, 1993.

\bibitem{norell2007}
U.~Norell.
\newblock Towards a practical programming language based on dependent type
  theory.
\newblock PhD thesis, Chalmers, 2007.

\bibitem{rfc6762}
S.~Cheshire and M.~Krochmal.
\newblock Multicast {DNS}.
\newblock RFC 6762, 2013.

\bibitem{rfc6763}
S.~Cheshire and M.~Krochmal.
\newblock {DNS}-Based Service Discovery.
\newblock RFC 6763, 2013.

\bibitem{rfc8305}
D.~Schinazi and T.~Pauly.
\newblock Happy Eyeballs Version 2: Better Connectivity Using Concurrency.
\newblock RFC 8305, 2017.

\bibitem{rfc8615}
M.~Nottingham.
\newblock Well-Known Uniform Resource Identifiers ({URIs}).
\newblock RFC 8615, 2019.

\end{thebibliography}

\end{document}
